\section{本章小结}
\label{sec:graph-traversal-summary}

本章的出发点是：第~\ref{chap:graph-objects}章定义的可达、分量、无环和同构，怎样在一张实际的图上算出来。回答的共同方式是沿邻接表有组织地访问顶点，用少量状态记住已经知道的东西，而不去枚举路径或双射。表~\ref{tab:graph-trav-summary}对照了本章各方法所维护的状态、输出和计算时间。

\begin{table*}[!htb]
  \centering
  \caption{本章方法维护的状态、输出与时间。时间均以邻接表存储计，$n$为顶点数，$m$为边数或弧数。}
  \label{tab:graph-trav-summary}
  \small
  \begin{tabular}{@{}llll@{}}
    \toprule
    方法 & 维护的状态 & 输出 & 时间 \\
    \midrule
    广度优先遍历 & 队列、距离、父指针 & 最少边数与一条最短路线 & $O(n+m)$ \\
    深度优先遍历 & 三色状态、栈帧、时刻 & 发现次序、完成次序、回边 & $O(n+m)$ \\
    分量编号 & 外层扫描、分量号 & 连通分量、生成森林 & $O(n+m)$ \\
    并查集 & 父指针树、集合大小 & 加边过程中的分量代表 & 每次$O(\log n)$ \\
    Kosaraju算法 & 完成次序、反向图 & 强连通分量、缩点图 & $O(n+m)$ \\
    Kahn算法 & 当前入度、零入度队列 & 拓扑序或有环报告 & $O(n+m)$ \\
    分隔集分块 & 屏蔽标记、分量号 & 余图分量、局部块 & $O(n+m)$ \\
    1-WL细化 & 各轮标记、签名字典 & 稳定划分、标记计数 & 每轮$O((n+m)\log n)$ \\
    \bottomrule
  \end{tabular}
\end{table*}

这些输出回答的是不同的问题。广度优先遍历的层数回答“离起点最少几步”，深度优先遍历的完成次序回答“谁的后续全部探完”，它在无环图上给出拓扑序，在一般有向图上把强连通分量按缩点图的先后分离出来。并查集的代表元回答“在已加入的边下是否连通”，1-WL的标记回答“局部邻域在计数意义下是否一样”。同一张图上，这些量可以同时计算，但不能互相替代：深度优先的父路径不是最短路线，发现次序不是拓扑序，并查集的父指针不是原图的边，1-WL计数相同也不是同构的证明。

判断一个算法的输出是否正确，标准是看它是否满足相应的结构条件，而不只是看算法是否停止。广度优先遍历的父路径应恰有$d(v)$条边；生成森林应有$n-c$条边且无环；拓扑序应使每条弧从前指向后；有环报告应附带一个可以逐弧核对的环；强连通分量在第二遍中不应越过已标记的分量；分块检验应报告余图实际有几个分量。各节的定理正是把这些条件与算法的状态联系起来的依据。

本章得到的可达区域、生成森林和依赖次序，都只描述了“能否连接、按什么次序”，没有比较不同路线或不同边集的代价：广度优先遍历只数边数，生成森林只保证连通。下一章在这些可行结构的基础上引入边上的数值，研究最短路、最小生成树、最大流、最小费用流和匹配；本章的并查集用于Kruskal算法，拓扑序用于有向无环图上的最短路递推。第~\ref{chap:graph-dynamic-programming}章则沿着第~\ref{sec:graph-separator-blocks}节的接口思想，把局部块上的计算组织起来。
